\documentclass[11pt,letterpaper]{article}

\usepackage[spanish,es-noshorthands]{babel}
\usepackage{fontspec}
\setmainfont{Source Serif Pro}
\usepackage{microtype}
\usepackage{geometry}
\usepackage{enumitem}
\usepackage{xcolor}
\usepackage{csquotes}
\usepackage{float}
\usepackage{tikz}
\usetikzlibrary{arrows.meta,positioning,fit,backgrounds,calc}
\usepackage{xurl}
\usepackage[hidelinks]{hyperref}
\usepackage[backend=biber,style=apa,sorting=nyt,maxcitenames=2]{biblatex}
\DeclareLanguageMapping{spanish}{spanish-apa}

\geometry{margin=2.3cm}
\hbadness=10000
\setlength{\emergencystretch}{1em}
\setlength{\parindent}{0pt}
\setlength{\parskip}{0.5em}
\setlist{nosep}
\addbibresource{referencias.bib}
\AtBeginBibliography{\sloppy\small}
\setlength{\bibitemsep}{0.3em}

% Colores: naranja para Stack Overflow, verde para la IA generativa,
% azul oscuro para flechas y rótulos y gris para los fondos.
\definecolor{cSO}{RGB}{200,96,20}
\definecolor{cIA}{RGB}{36,122,82}
\definecolor{cEje}{RGB}{26,58,102}
\definecolor{cGris}{RGB}{120,116,110}

\tikzset{
  etapa/.style={draw=#1, fill=#1!8, very thick, rounded corners=4pt,
    text width=4.4cm, align=center, minimum height=1.35cm, font=\small},
  flujo/.style={-{Stealth[length=7pt]}, very thick, cEje},
}

\begin{document}

% --- Portada -----------------------------------------------------------------
\begin{titlepage}
\centering
\vspace*{3.2cm}
{\large\scshape Gestión del Desarrollo Tecnológico:\\
Estrategias y Análisis de Portfolio\par}
\vspace{0.5cm}
{\normalsize Módulo 8. Foro 8.1: Ejemplos de tecnologías abandonadas\par}
\vspace{2.6cm}
{\color{cEje}\rule{0.55\textwidth}{0.8pt}\par}
\vspace{0.7cm}
{\LARGE\bfseries Stack Overflow y la IA generativa\par}
\vspace{0.5cm}
{\large Cuando la fuente de entrenamiento\\
se convierte en víctima\par}
\vspace{0.7cm}
{\color{cEje}\rule{0.55\textwidth}{0.8pt}\par}
\vfill
{\large Carlos Luis Rojas Aragonés\par}
\vspace{0.3cm}
{\normalsize Chief Technology Officer Program\par}
\vspace{0.3cm}
{\normalsize 26 de octubre de 2026\par}
\vspace*{1.5cm}
\end{titlepage}

% --- Cuerpo ------------------------------------------------------------------
\section*{1. Del sitio de referencia a la caída}

Durante más de una década, Stack Overflow fue el sitio de referencia para
resolver dudas de programación. Su modelo era sencillo: una comunidad que
respondía gratis y un buscador que posicionaba esas respuestas. En 2021, Prosus
la compró por 1800 millones de dólares \parencite{stackoverflow2021}. Un año y
medio después llegó ChatGPT.

El efecto fue inmediato. \Textcite{delriochanona2024} estimaron que la actividad
en Stack Overflow cayó cerca de un 25\,\% en los seis meses siguientes al
lanzamiento de ChatGPT, frente a plataformas equivalentes en Rusia y China,
donde la herramienta no estaba disponible. En octubre de 2023, la empresa
despidió al 28\,\% de su plantilla. La dirección lo atribuyó a la búsqueda de
rentabilidad, pero coincidió con esa caída \parencite{chandrasekar2023}.

\section*{2. Una disrupción de manual}

El caso encaja con el dilema del innovador \parencite{christensen1997}. El
asistente de IA es \enquote{peor} en precisión: alucina y no cita fuentes. Sin
embargo, responde al instante, no juzga la pregunta y se adapta al código
exacto del desarrollador. Para la mayoría de las dudas cotidianas, eso basta. En
la encuesta de la propia Stack Overflow, el 76\,\% de los desarrolladores ya
usaba o planeaba usar herramientas de IA \parencite{stackoverflow2024encuesta}.

\section*{3. La ironía de fondo}

Los modelos aprendieron del contenido de Stack Overflow, y ahora reducen las
preguntas nuevas que alimentarían a los modelos futuros
(Figura~\ref{fig:circulo}). La empresa respondió licenciando sus datos a Google
y OpenAI \parencite{stackoverflow2024google,stackoverflow2024openai}. Es decir, pasó de ser una
comunidad a ser un proveedor de datos.

\begin{figure}[H]
\centering
\begin{tikzpicture}[node distance=1.5cm and 1.6cm]
  \node[etapa=cSO] (resp) {Respuestas públicas\\en Stack Overflow};
  \node[etapa=cIA, right=of resp] (entr) {Entrenamiento de\\los modelos de lenguaje};
  \node[etapa=cIA, below=of entr] (asis) {Los desarrolladores\\preguntan al asistente};
  \node[etapa=cSO, below=of resp] (menos) {Menos preguntas nuevas\\(cerca de $-$25\,\% en seis meses)};
  \draw[flujo] (resp) -- (entr);
  \draw[flujo] (entr) -- (asis);
  \draw[flujo] (asis) -- (menos);
  \draw[flujo] (menos) -- node[left, font=\small, text=cEje, align=right]
    {menos\\contenido} (resp);
  \node[font=\small\itshape, text=cGris] at ($(resp)!0.5!(asis)$) {círculo vicioso};
\end{tikzpicture}
\caption{El contenido de la comunidad alimenta a los modelos que luego la
sustituyen. Elaboración propia con base en \Textcite{delriochanona2024}.}
\label{fig:circulo}
\end{figure}

Les dejo una pregunta: si las comunidades abiertas dejan de generar
conocimiento, ¿de dónde aprenderán los próximos modelos?

\printbibliography[title={Referencias}]

\end{document}
